\begin{table}[htbp]
\centering
\caption{Summary Statistics}
\label{tab:summary}
\begin{tabular}{lcccc}
\toprule
 & \multicolumn{1}{c}{Full Sample} & \multicolumn{3}{c}{By Wage Ratio Tercile} \\
\cmidrule(lr){2-2} \cmidrule(lr){3-5}
 & & Low & Medium & High \\
\midrule
\multicolumn{5}{l}{\textit{Panel A: Treatment Variable}} \\
Personal care aide wage (\$/hr) & 9.66 & 8.87 & 9.14 & 10.98 \\
Competing sector wage (\$/hr) & 14.23 & 16.22 & 13.29 & 13.18 \\
Wage competitiveness ratio & 0.700 & 0.578 & 0.687 & 0.833 \\
[4pt]
\multicolumn{5}{l}{\textit{Panel B: Outcomes (state-month means)}} \\
Active HCBS providers & 706 & 906 & 538 & 674 \\
Beneficiaries served & 176,603 & 264,289 & 110,972 & 154,548 \\
Monthly spending (\$M) & 96.9 & 114.2 & 54.8 & 121.7 \\
[4pt]
\multicolumn{5}{l}{\textit{Panel C: Sample}} \\
States & 51 & 17 & 17 & 17 \\
Months & 83 & & & \\
State-month observations & 4,233 & & & \\
\bottomrule
\end{tabular}
\begin{tablenotes}[flushleft]\small
\item \textit{Notes:} Wage competitiveness ratio = BLS QCEW average hourly wage in NAICS 624120 (Services for the Elderly and Persons with Disabilities) divided by simple average of hourly wages in three competing industries (grocery stores, limited-service restaurants, general warehousing). All wages are 2019 annual averages, private sector. Tercile cutoffs based on 2019 ratio distribution. HCBS providers defined as billing NPIs with T-code or S-code activity in T-MSIS.
\end{tablenotes}
\end{table}
